\section{The sharp hypersurface Sobolev inequality}
\label{sec:sobolev}

We now turn the comparison maps into a sharp inequality on a possibly
incomplete hypersurface. The proof first rules out a Dirichlet quotient
below the Euclidean constant; a point cutoff then permits tests supported
on an entire compact component.

Let $b\in\{0,1\}$, and let $M^n$ be complete and simply connected,
with $\Sec_g\le-b^2$. Put $m=n-1\ge3$ and
\[
 q=\frac{2m}{m-2},\qquad c=\frac4{m-2}=q-2,\qquad
 \alpha=\frac{m-2}{2},\qquad Y_m=m s_m^{2/m}.
\]
Let $\Sigma^m$ be a locally $C^{1,1}$ immersed hypersurface in $M$,
with its induced metric and area measure $d\sigma$. It need not be
complete and may have boundary. Compactness is understood in the
manifold topology of $\Sigma$, and $\Sigma^\circ$ denotes its interior.
Ambient functions and maps on $\Sigma$ are composed with the immersion.
For a local unit normal $N$, write $mh=\divg_\Sigma N$ almost everywhere.
The square $h^2$ is independent of the choice of normal. Define
\begin{equation}
\label{eq:sobolev-form}
 Q_\Sigma(v)=\int_\Sigma
 \left(c|\nabla_\Sigma v|^2+m(h^2-b^2)v^2\right)\,d\sigma.
\end{equation}

For open $U\Subset\Sigma^\circ$, let $H^1_0(U)$ be the closure of
compactly supported Lipschitz functions in the induced $H^1$ norm;
no regularity of $\partial U$ is assumed.

\begin{theorem}[Hypersurface Sobolev inequality]
\label{thm:extrinsic-sobolev}
For every Lipschitz function $v$ with compact support in $\Sigma^\circ$,
\begin{equation}
\label{eq:extrinsic-sobolev}
 Q_\Sigma(v)\ge
 Y_m\left(\int_\Sigma |v|^q\,d\sigma\right)^{2/q}.
\end{equation}
The inequality also holds for the zero extension of every
$v\in H^1_0(U)$, where $U\Subset\Sigma^\circ$ is open.
\end{theorem}

The constant is sharp: on a geodesic sphere in the space form of
curvature $-b^2$, the constant function gives equality.

By this definition, zero extension satisfies
\begin{equation}
\label{eq:sobolev-zero-extension}
 E_0v\in H^1(\Sigma^\circ),\qquad
 \nabla_\Sigma(E_0v)=\mathbf1_U\nabla_\Sigma v
 \quad\text{almost everywhere}.
\end{equation}
Both statements hold for the approximating functions and pass to their
$H^1$ limits, including when $\partial U$ has positive measure.
In $C^{1,1}$ charts the induced metric is locally Lipschitz and uniformly
positive definite on compact subsets, and $h$ is locally essentially
bounded. Thus the potential in~\eqref{eq:sobolev-form} is bounded
near $\bar U$ and $Q_\Sigma$ is continuous on $H^1_0(U)$.

\subsection{A geometric square}

The comparison maps provide the following lower bound before any
minimization takes place.

\begin{lemma}
\label{lem:square-estimate}
Fix $p\in M$ and $a>0$, and let $\Phi=\Phi_{p,a}$ and
$f=e^u=f_{p,a}$ be the map and functions of
Lemma~\ref{lem:comparison-maps}. For every compactly supported
Lipschitz function $v$ in $\Sigma^\circ$,
\begin{equation}
\label{eq:square-estimate}
\begin{aligned}
 Q_\Sigma(v)
 &\ge\int_\Sigma\left(
 mf^2v^2+m(h+Nu)^2v^2
 +c|\nabla_\Sigma v-\alpha v\nabla_\Sigma u|^2\right)\,d\sigma\\
 &\ge\int_\Sigma\left(
 v^2|d(\Phi|_\Sigma)|^2
 +c|\nabla_\Sigma v-\alpha v\nabla_\Sigma u|^2\right)\,d\sigma.
\end{aligned}
\end{equation}
The same inequalities hold on $H^1_0(U)$ for every
$U\Subset\Sigma^\circ$. In particular, $Q_\Sigma$ is nonnegative
on all these functions.
\end{lemma}

\begin{proof}
Write $\nabla=\nabla_\Sigma$ and $\beta=Nu$. The restriction formula is
\begin{equation}
\label{eq:sobolev-restriction-laplacian}
 \Delta_\Sigma u=\tr_{T\Sigma}\Hess_g u-mh\beta.
\end{equation}
This holds both almost everywhere and weakly: in a $C^{1,1}$
parametrization, the induced divergence expression has Lipschitz
coefficients, and the restricted smooth ambient function has Lipschitz
first derivatives. The ordinary chain rule therefore gives the same
formula as the distributional divergence.

Taking the tangential trace in Lemma~\ref{lem:comparison-maps}, and
using $|\nabla_g u|^2=|\nabla u|^2+\beta^2$, gives
\[
 2\Delta_\Sigma u
 \le -(m-2)|\nabla u|^2-m(b^2+f^2)-m\beta^2-2mh\beta.
\]
Equivalently,
\[
 m(h^2-b^2)\ge
 mf^2+m(h+\beta)^2+(m-2)|\nabla u|^2+2\Delta_\Sigma u.
\]
Multiply by $v^2$ and integrate by parts. Since
$c\alpha=2$ and $c\alpha^2=m-2$, the gradient terms become
\[
 c|\nabla v|^2-4v\langle\nabla v,\nabla u\rangle
 +(m-2)v^2|\nabla u|^2
 =c|\nabla v-\alpha v\nabla u|^2.
\]
This proves the first inequality. The squared differential norm is the
sum over an orthonormal tangent frame, so
$|d(\Phi|_\Sigma)|^2\le mf^2$ by
Lemma~\ref{lem:comparison-maps}, proving the second.
Reversing a local normal reverses both $h$ and $\beta$; all expressions
are therefore globally defined. Finally, all coefficients are bounded
near $\bar U$, so $H^1_0(U)$ approximation proves the last assertion.
\end{proof}

\subsection{Compactness below the Euclidean constant}

We need only local compactness on the hypersurface. The below-threshold
argument follows the critical-quotient application of
Br\'ezis--Lieb~\cite[Section~4(B)]{BrezisLieb1983}; we give the details
for the present metric and potential. The critical concentration
threshold is the same for both curvature bounds.

\begin{lemma}
\label{lem:local-sobolev-threshold}
For every open $U\Subset\Sigma^\circ$ and $\epsilon>0$, there is
$C_{\epsilon,U}<\infty$ such that
\begin{equation}
\label{eq:local-sobolev-threshold}
 (Y_m-\epsilon)\|v\|_q^2
 \le c\int_\Sigma|\nabla_\Sigma v|^2\,d\sigma
       +C_{\epsilon,U}\int_\Sigma v^2\,d\sigma,
 \qquad v\in H^1_0(U).
\end{equation}
The inclusion $H^1_0(U)\hookrightarrow L^2(U)$ is compact.
\end{lemma}

\begin{proof}
The sharp Euclidean Sobolev inequality of Aubin and
Talenti~\cite{Aubin1976,Talenti1976}, in our normalization, is
\begin{equation}
\label{eq:talenti-normalization}
 c\int_{\mathbb R^m}|D\varphi|^2\,dx
 \ge Y_m\left(\int_{\mathbb R^m}|\varphi|^q\,dx\right)^{2/q},
 \qquad \varphi\in\Lip_c(\mathbb R^m).
\end{equation}
Before multiplying by $c$, the sharp constant is
$m(m-2)|\mathbb S^m|^{2/m}/4$.
Normalize the metric at a chart center. In a sufficiently small chart,
its matrix $G$ satisfies
$(1-\delta)\mathrm{Id}\le G\le(1+\delta)\mathrm{Id}$.
Comparing inverse metrics and volume densities in
\eqref{eq:talenti-normalization} gives
\begin{equation}
\label{eq:chart-sobolev-distortion}
 c\int_\Sigma|\nabla_\Sigma\varphi|^2\,d\sigma
 \ge Y_m\left(\frac{1-\delta}{1+\delta}\right)^{m/2}
          \|\varphi\|_q^2
\end{equation}
for chart-supported $\varphi$. The gradient integral is at least
$(1-\delta)^{m/2}(1+\delta)^{-1}$ times its Euclidean value,
whereas the squared critical norm is at most
$(1+\delta)^{m/q}$ times its Euclidean value, and $m/q=(m-2)/2$.

Choose $\delta$ so that the coefficient in
\eqref{eq:chart-sobolev-distortion} is at least $Y_m-\epsilon$.
Cover $\bar U$ by finitely many such charts and choose chart-supported
Lipschitz functions $\chi_i$ with $\sum_i\chi_i^2=1$ near $\bar U$.
They are obtained by normalizing a finite family of cutoffs by the
square root of the sum of their squares near $\bar U$, and then
cutting off outside that neighborhood. The triangle inequality in
$L^{q/2}$ and the product rule give
\begin{align}
\label{eq:squared-partition-norm}
 \|v\|_q^2
 &=\left\|\sum_i(\chi_i v)^2\right\|_{q/2}
 \le\sum_i\|\chi_i v\|_q^2,\\
\label{eq:squared-partition-energy}
 \sum_i|\nabla_\Sigma(\chi_i v)|^2
 &=|\nabla_\Sigma v|^2
        +v^2\sum_i|\nabla_\Sigma\chi_i|^2.
\end{align}
Summing~\eqref{eq:chart-sobolev-distortion} proves
\eqref{eq:local-sobolev-threshold} first for compactly supported
Lipschitz functions. Taking, for example, $\epsilon=Y_m/2$ gives
the continuous inclusion into $L^q$, so approximation proves it
for every $v\in H^1_0(U)$.

For compactness, extend a bounded $H^1_0(U)$ sequence by zero.
Each $\chi_i v_j$ is a bounded Euclidean $H^1$ sequence in its chart,
supported in a fixed compact subset. Rellich compactness gives an
$L^2$-convergent subsequence in each of the finitely many charts.
The identity $v_j=\sum_i\chi_i(\chi_i v_j)$ then gives convergence
in $L^2(U)$. No boundary regularity of $U$ is used.
\end{proof}

\begin{lemma}[Attainment below the threshold]
\label{lem:subthreshold-minimizer}
For nonempty open $U\Subset\Sigma^\circ$, put
\[
 Y(U)=\inf\{Q_\Sigma(v):v\in H^1_0(U),\ \|v\|_q=1\}.
\]
If $Y(U)<Y_m$, the infimum is attained.
\end{lemma}

\begin{proof}
Write $Y=Y(U)\ge0$, the last inequality following from
Lemma~\ref{lem:square-estimate}. For a normalized minimizing sequence,
H\"older's inequality and $h^2-b^2\ge-b^2$ give
\begin{equation}
\label{eq:minimizing-sequence-bound}
 \|v_j\|_2^2\le\sigma(U)^{2/m},\qquad
 c\|\nabla_\Sigma v_j\|_2^2
 \le Q_\Sigma(v_j)+mb^2\sigma(U)^{2/m}.
\end{equation}
Thus, after a subsequence,
\[
 v_j\rightharpoonup\psi\text{ in }H^1_0(U),\qquad
 v_j\longrightarrow\psi\text{ in }L^2(U)
 \text{ and almost everywhere}.
\]
Set $t=\int_U|\psi|^q\,d\sigma\in[0,1]$ and $z_j=v_j-\psi$.
The Br\'ezis--Lieb lemma~\cite[Theorem~1]{BrezisLieb1983} gives
\begin{equation}
\label{eq:brezis-lieb-splitting}
 \int_U|z_j|^q\,d\sigma\longrightarrow1-t.
\end{equation}
Weak convergence splits the quadratic form:
$Q_\Sigma(v_j)=Q_\Sigma(\psi)+Q_\Sigma(z_j)+o(1)$.
The bounded potential and~\eqref{eq:local-sobolev-threshold} give
\[
 Q_\Sigma(z_j)\ge
 (Y_m-\epsilon)\|z_j\|_q^2
 -(C_{\epsilon,U}+mb^2)\|z_j\|_2^2.
\]
Letting $j\to\infty$ and then $\epsilon\downarrow0$, we obtain
\begin{equation}
\label{eq:critical-mass-inequality}
 Y\ge Q_\Sigma(\psi)+Y_m(1-t)^{2/q}
 \ge Yt^{2/q}+Y_m(1-t)^{2/q}.
\end{equation}
If $t=0$, this contradicts $Y<Y_m$. If $Y=0$, it forces $t=1$.
If $Y>0$ and $0<t<1$, strict concavity gives
$t^{2/q}+(1-t)^{2/q}>1$, and the right side is strictly greater
than $Y$. Thus $t=1$ in every case. Lower semicontinuity of the
gradient energy and strong $L^2$ convergence of the potential term
give $Q_\Sigma(\psi)\le Y$; normalization gives the reverse inequality.
\end{proof}

\subsection{Balanced variations}

A quotient below the sharp threshold would now have a normalized
minimizer $\psi$. After balancing its critical mass with a comparison
map $\Phi$, the spherical-coordinate variations will give
$\int_\Sigma\psi^2|d(\Phi|_\Sigma)|^2\,d\sigma\ge Q_\Sigma(\psi)$.
The geometric square gives the reverse inequality with an additional
nonnegative gradient square. The two bounds force that square to vanish,
and zero extension supplies the contradiction.

\begin{proof}[Proof of Theorem~\ref{thm:extrinsic-sobolev}]
First let $\Sigma$ be connected and boundaryless, and fix a nonempty
open $U\Subset\Sigma$ with $\Sigma\setminus\bar U\ne\varnothing$.
Suppose $Y(U)<Y_m$, and take the normalized minimizer $\psi$ supplied
by Lemma~\ref{lem:subthreshold-minimizer}. Write $Y=Q_\Sigma(\psi)$.

Push the probability measure $|\psi|^q\,d\sigma$, extended by zero,
to $M$ by the immersion. Its support is compact and it has no atoms.
Indeed, an immersion is locally an embedding, so a fiber in the compact
support meets each member of a finite embedding-chart cover in at most
one point; every such fiber has zero area measure.
Lemma~\ref{lem:balancing} supplies fixed parameters $p,a$ such that
the coordinates $w_1,\ldots,w_{n+1}$ of $\Phi_{p,a}|_\Sigma$ satisfy
\begin{equation}
\label{eq:balanced-minimizer}
 \int_\Sigma|\psi|^q w_j\,d\sigma=0,\qquad
 \sum_{j=1}^{n+1}w_j^2=1.
\end{equation}
These coordinates and $u=u_{p,a}$ have bounded first derivatives
near $\bar U$. Multiplication by $w_j$ or $w_j^2$ preserves $H^1_0(U)$.

Let $B_Q$ be the symmetric bilinear form associated with $Q_\Sigma$.
For a bounded Lipschitz multiplier $\zeta$ near $\bar U$,
differentiate the quotient along $\psi(1+s\zeta)$ to obtain
\begin{equation}
\label{eq:minimizer-first-variation}
 B_Q(\psi,\psi\zeta)
 =Y\int_\Sigma|\psi|^q\zeta\,d\sigma.
\end{equation}
For small $|s|$, the factor $1+s\zeta$ is bounded away from zero,
and differentiated mass integrands are dominated by a constant times
$|\psi|^q$. No regularity or positivity of the minimizer is needed.

Fix a coordinate $w=w_j$ and put
$B_w=\int_\Sigma|\psi|^q w^2\,d\sigma$. Since $|w|\le1$ and
its weighted mean is zero,
\[
 \|\psi(1+sw)\|_q^2=1+(q-1)B_ws^2+o(s^2).
\]
The numerator has zero linear term by
\eqref{eq:minimizer-first-variation}; its second variation at the
minimum therefore gives
\begin{equation}
\label{eq:minimizer-second-variation}
 Q_\Sigma(\psi w)\ge(q-1)YB_w.
\end{equation}
The product rule and the first variation with $\zeta=w^2$ give
\begin{equation}
\label{eq:coordinate-product-identity}
\begin{aligned}
 Q_\Sigma(\psi w)
 &=B_Q(\psi,\psi w^2)
       +c\int_\Sigma\psi^2|\nabla_\Sigma w|^2\,d\sigma\\
 &=YB_w+c\int_\Sigma\psi^2|\nabla_\Sigma w|^2\,d\sigma.
\end{aligned}
\end{equation}
Subtract and sum over the coordinates. Since $c=q-2$ and
$\sum_jw_j^2=1$, the result is
\begin{equation}
\label{eq:balanced-map-energy}
 \int_\Sigma\psi^2|d(\Phi_{p,a}|_\Sigma)|^2\,d\sigma\ge Y.
\end{equation}
Combining this with Lemma~\ref{lem:square-estimate} forces
\begin{equation}
\label{eq:weighted-minimizer-gradient}
 \nabla_\Sigma\psi=\alpha\psi\nabla_\Sigma u
 \quad\text{almost everywhere on }U.
\end{equation}

Choose a compactly supported Lipschitz cutoff $\chi$ on $\Sigma$
equal to one near $\bar U$. By~\eqref{eq:sobolev-zero-extension},
\[
 F=(\chi e^{-\alpha u})E_0\psi\in H^1(\Sigma)
\]
is the zero extension of $e^{-\alpha u}\psi$, and
\eqref{eq:weighted-minimizer-gradient} gives $\nabla_\Sigma F=0$
almost everywhere on all of $\Sigma$. The local Poincar\'e inequality
in coordinate balls makes $F$ constant almost everywhere on each ball.
Overlapping balls have the same constant, so connectedness makes $F$
constant on $\Sigma$. It vanishes on the nonempty open set
$\Sigma\setminus\bar U$ and therefore vanishes everywhere,
contradicting $\|\psi\|_q=1$. Thus $Y(U)\ge Y_m$.

Every proper compact support in a connected boundaryless $\Sigma$
is contained in such a $U$: choose a coordinate ball outside the
support and finitely many relatively compact neighborhoods covering
the support that miss a smaller ball. This proves the inequality
unless the support is the entire compact component.
In that case, let $D_\delta$ be the image of a Euclidean ball of radius
$\delta$ in a fixed chart about a point. Remove that point with a
cutoff $\eta_\delta$ equal to zero in $D_\delta$, equal to one outside
$D_{2\delta}$, and satisfying
$|\nabla_\Sigma\eta_\delta|\le C/\delta$.
For Lipschitz $v$, metric comparability gives
\[
 \int_\Sigma|\nabla_\Sigma((1-\eta_\delta)v)|^2\,d\sigma
 \le 2\int_{D_{2\delta}}|\nabla_\Sigma v|^2\,d\sigma
       +C\|v\|_\infty^2\delta^{m-2}\longrightarrow0.
\]
Also $\eta_\delta v\to v$ in $L^2$ and $L^q$. The potential is
bounded on the compact component, so the inequality for
$\eta_\delta v$ passes to $v$. This is where $m>2$ is needed.

Finally, pass to $\Sigma^\circ$ if $\Sigma$ has boundary.
A compact support meets only finitely many connected components,
because the components are open. If
$a_j=\int_{\Sigma_j}|v|^q\,d\sigma$, the connected result gives
\[
 Q_\Sigma(v)\ge Y_m\sum_j a_j^{2/q}
 \ge Y_m\left(\sum_j a_j\right)^{2/q}.
\]
For $v\in H^1_0(U)$, use its defining approximation:
Lemma~\ref{lem:local-sobolev-threshold} gives convergence in $L^q$,
and boundedness of the potential near $\bar U$ gives convergence
of the quadratic form. This proves the final assertion.
\end{proof}
